% One-time imported draft. Editable native source; no Markdown regeneration.
\section{Derive the framebuffer
addresses}\label{003-eight-byte-sprite-p02-derive-the-framebuffer-addresses}

A memory address identifies a byte. The destination uses four bytes for
each RGB word, so pixel coordinates must be scaled before they can be
used as byte displacements. With sixteen columns in each destination
row, the word index of column x in row y is \texttt{16\ \ensuremath{\times}\ y\ +\ x}.
Multiplying that index by four gives the byte displacement from the
frame base.

The original image uses x from zero through seven. The mirrored image
uses x from eight through fifteen. The same row stride, 64 bytes,
applies to both. Advancing a compact source row requires one byte, while
advancing a framebuffer row requires 64. Confusing these units can write
a recognizable fragment in the wrong part of memory.

The greatest allowed index is \texttt{16\ \ensuremath{\times}\ 7\ +\ 15\ =\ 127}. Its byte
displacement is 508. That word ends at displacement 511, immediately
before the canary at 512. This calculation establishes the bounds that
the pointer loops must respect.
